{\color{ink}\bfseries 例 3-1 \(\cdot \) 合订本12 \(\cdot \) 图1}\par \textbf{栈顶在上；\(\uparrow \)入栈、\(\downarrow \)出栈}\par\vspace{2bp}\begin{tikzpicture}[x=1bp,y=-1bp,>=Stealth,font=\fontsize{8.5}{10.7}\selectfont,every node/.style={inner sep=1bp},line width=.45bp]
\node[anchor=north west,align=left] at (0,0){输出 b};
\node[anchor=north west,align=left,text width=52bp] at (0,15){\(\uparrow \)a \(\uparrow \)b \(\downarrow \)b};
\node[anchor=north west,align=left] at (62,0){峰值};
\node[anchor=north west,align=left] at (101,0){结束};
\draw (64,15) -- (64,43) -- (87,43) -- (87,15);
\node[anchor=north west,align=left] at (72,16){b};
\draw (64,29) -- (87,29);
\node[anchor=north west,align=left] at (72,30){a};
\draw (64,43) -- (87,43);
\draw (105,15) -- (105,29) -- (128,29) -- (128,15);
\node[anchor=north west,align=left] at (113,16){a};
\draw (105,29) -- (128,29);
\node[anchor=north west,align=left] at (0,76){输出 d,\allowbreak{}c};
\node[anchor=north west,align=left,text width=52bp] at (0,91){\(\uparrow \)c \(\uparrow \)d \(\downarrow \)d \(\downarrow \)c};
\node[anchor=north west,align=left] at (62,76){峰值};
\node[anchor=north west,align=left] at (101,76){结束};
\draw (64,91) -- (64,133) -- (87,133) -- (87,91);
\node[anchor=north west,align=left] at (72,92){d};
\draw (64,105) -- (87,105);
\node[anchor=north west,align=left] at (72,106){c};
\draw (64,119) -- (87,119);
\node[anchor=north west,align=left] at (72,120){a};
\draw (64,133) -- (87,133);
\draw (105,91) -- (105,105) -- (128,105) -- (128,91);
\node[anchor=north west,align=left] at (113,92){a};
\draw (105,105) -- (128,105);
\node[anchor=north west,align=left] at (0,152){输出 f,\allowbreak{}e};
\node[anchor=north west,align=left,text width=52bp] at (0,167){\(\uparrow \)e \(\uparrow \)f \(\downarrow \)f \(\downarrow \)e};
\node[anchor=north west,align=left] at (62,152){峰值};
\node[anchor=north west,align=left] at (101,152){结束};
\draw (64,167) -- (64,209) -- (87,209) -- (87,167);
\node[anchor=north west,align=left] at (72,168){f};
\draw (64,181) -- (87,181);
\node[anchor=north west,align=left] at (72,182){e};
\draw (64,195) -- (87,195);
\node[anchor=north west,align=left] at (72,196){a};
\draw (64,209) -- (87,209);
\draw (105,167) -- (105,181) -- (128,181) -- (128,167);
\node[anchor=north west,align=left] at (113,168){a};
\draw (105,181) -- (128,181);
\node[anchor=north west,align=left] at (0,228){输出 a,\allowbreak{}g};
\node[anchor=north west,align=left,text width=52bp] at (0,243){\(\downarrow \)a \(\uparrow \)g \(\downarrow \)g};
\node[anchor=north west,align=left] at (62,228){峰值};
\node[anchor=north west,align=left] at (101,228){结束};
\draw (64,243) -- (64,257) -- (87,257) -- (87,243);
\node[anchor=north west,align=left] at (72,244){a};
\draw (64,257) -- (87,257);
\draw (105,243) -- (105,257) -- (128,257) -- (128,243);
\node[anchor=north west,align=left] at (112,244){空};
\end{tikzpicture}\par 最大深度3，最小容量3（C）；每段结束仅余a，最终为空。
